% An inode with direct, single, double and triple indirect pointers.
% Usage: % An inode with direct, single, double and triple indirect pointers.
% Usage: % An inode with direct, single, double and triple indirect pointers.
% Usage: % An inode with direct, single, double and triple indirect pointers.
% Usage: \input{figures/l11-inode}   (width ~116 mm)
\begin{tikzpicture}[x=1mm, y=1mm, font=\scriptsize\sffamily,
  >={Stealth[length=1.4mm]},
  data/.style={draw, fill=white, minimum width=10mm, minimum height=4.5mm,
               inner sep=0.5pt, font=\scriptsize\sffamily},
  lab/.style={font=\scriptsize\sffamily, align=center}]
  % \ptrblk{name}{x}{y}: block of pointers
  \def\ptrblk#1#2#3{%
    \draw[fill=ln-box] (#2-5,#3-4) rectangle (#2+5,#3+4);
    \foreach \k in {-2,0,2} \draw[ln-rule] (#2-5,#3+\k) -- (#2+5,#3+\k);
    \coordinate (#1w) at (#2-5,#3);
    \coordinate (#1e) at (#2+5,#3+3);
    \coordinate (#1s) at (#2+5,#3-3);}
  % inode
  \node[font=\footnotesize\sffamily\bfseries, anchor=south] at (19,48.3) {\iflangja{iノード}{inode}};
  \draw[fill=ln-accent!10] (4,36) rectangle (34,48);
  \node[lab] at (19,42) {\iflangja{メタデータ\\（モード，所有者，\\タイムスタンプ，サイズ）}{metadata\\(mode, owners,\\timestamps, size, \ldots)}};
  \def\irow#1#2{\draw[fill=white] (4,#1-2.5) rectangle (34,#1+2.5);
    \node[lab] at (19,#1) {#2};}
  \irow{33.5}{\iflangja{直接ポインタ 0}{direct 0}}
  \irow{28.5}{\ldots}
  \irow{23.5}{\iflangja{直接ポインタ 11}{direct 11}}
  \irow{18.5}{\iflangja{間接ポインタ}{single indirect}}
  \irow{13.5}{\iflangja{二重間接ポインタ}{double indirect}}
  \irow{8.5}{\iflangja{三重間接ポインタ}{triple indirect}}
  % direct
  \node[data] (a) at (111,33.5) {\iflangja{データ}{data}};
  \node[data] (b) at (111,23.5) {\iflangja{データ}{data}};
  \draw[->] (31,33.5) -- (a.west);
  \draw[->] (31,23.5) -- (b.west);
  % single indirect
  \ptrblk{p}{53}{15}
  \node[data] (c) at (111,15) {\iflangja{データ}{data}};
  \draw[->] (31,18.5) -- (pw);
  \draw[->] (pe) -- (c.west);
  % double indirect
  \ptrblk{q}{53}{0}
  \ptrblk{r}{77}{0}
  \node[data] (d) at (111,0) {\iflangja{データ}{data}};
  \draw[->] (31,13.5) -- (qw);
  \draw[->] (qe) -- (rw);
  \draw[->] (re) -- (d.west);
  % triple indirect
  \ptrblk{s}{53}{-15}
  \ptrblk{t}{71}{-15}
  \ptrblk{u}{89}{-15}
  \node[data] (e) at (111,-15) {\iflangja{データ}{data}};
  \draw[->] (31,8.5) -- (sw);
  \draw[->] (se) -- (tw);
  \draw[->] (te) -- (uw);
  \draw[->] (ue) -- (e.west);
  % legend
  \draw[fill=ln-box] (4,-27) rectangle (10,-23);
  \node[lab, anchor=west] at (11,-25) {\iflangja{ポインタのブロック}{block of pointers}};
  \draw[fill=white] (44,-27) rectangle (50,-23);
  \node[lab, anchor=west] at (51,-25) {\iflangja{データブロック}{data block}};
\end{tikzpicture}
   (width ~116 mm)
\begin{tikzpicture}[x=1mm, y=1mm, font=\scriptsize\sffamily,
  >={Stealth[length=1.4mm]},
  data/.style={draw, fill=white, minimum width=10mm, minimum height=4.5mm,
               inner sep=0.5pt, font=\scriptsize\sffamily},
  lab/.style={font=\scriptsize\sffamily, align=center}]
  % \ptrblk{name}{x}{y}: block of pointers
  \def\ptrblk#1#2#3{%
    \draw[fill=ln-box] (#2-5,#3-4) rectangle (#2+5,#3+4);
    \foreach \k in {-2,0,2} \draw[ln-rule] (#2-5,#3+\k) -- (#2+5,#3+\k);
    \coordinate (#1w) at (#2-5,#3);
    \coordinate (#1e) at (#2+5,#3+3);
    \coordinate (#1s) at (#2+5,#3-3);}
  % inode
  \node[font=\footnotesize\sffamily\bfseries, anchor=south] at (19,48.3) {\iflangja{iノード}{inode}};
  \draw[fill=ln-accent!10] (4,36) rectangle (34,48);
  \node[lab] at (19,42) {\iflangja{メタデータ\\（モード，所有者，\\タイムスタンプ，サイズ）}{metadata\\(mode, owners,\\timestamps, size, \ldots)}};
  \def\irow#1#2{\draw[fill=white] (4,#1-2.5) rectangle (34,#1+2.5);
    \node[lab] at (19,#1) {#2};}
  \irow{33.5}{\iflangja{直接ポインタ 0}{direct 0}}
  \irow{28.5}{\ldots}
  \irow{23.5}{\iflangja{直接ポインタ 11}{direct 11}}
  \irow{18.5}{\iflangja{間接ポインタ}{single indirect}}
  \irow{13.5}{\iflangja{二重間接ポインタ}{double indirect}}
  \irow{8.5}{\iflangja{三重間接ポインタ}{triple indirect}}
  % direct
  \node[data] (a) at (111,33.5) {\iflangja{データ}{data}};
  \node[data] (b) at (111,23.5) {\iflangja{データ}{data}};
  \draw[->] (31,33.5) -- (a.west);
  \draw[->] (31,23.5) -- (b.west);
  % single indirect
  \ptrblk{p}{53}{15}
  \node[data] (c) at (111,15) {\iflangja{データ}{data}};
  \draw[->] (31,18.5) -- (pw);
  \draw[->] (pe) -- (c.west);
  % double indirect
  \ptrblk{q}{53}{0}
  \ptrblk{r}{77}{0}
  \node[data] (d) at (111,0) {\iflangja{データ}{data}};
  \draw[->] (31,13.5) -- (qw);
  \draw[->] (qe) -- (rw);
  \draw[->] (re) -- (d.west);
  % triple indirect
  \ptrblk{s}{53}{-15}
  \ptrblk{t}{71}{-15}
  \ptrblk{u}{89}{-15}
  \node[data] (e) at (111,-15) {\iflangja{データ}{data}};
  \draw[->] (31,8.5) -- (sw);
  \draw[->] (se) -- (tw);
  \draw[->] (te) -- (uw);
  \draw[->] (ue) -- (e.west);
  % legend
  \draw[fill=ln-box] (4,-27) rectangle (10,-23);
  \node[lab, anchor=west] at (11,-25) {\iflangja{ポインタのブロック}{block of pointers}};
  \draw[fill=white] (44,-27) rectangle (50,-23);
  \node[lab, anchor=west] at (51,-25) {\iflangja{データブロック}{data block}};
\end{tikzpicture}
   (width ~116 mm)
\begin{tikzpicture}[x=1mm, y=1mm, font=\scriptsize\sffamily,
  >={Stealth[length=1.4mm]},
  data/.style={draw, fill=white, minimum width=10mm, minimum height=4.5mm,
               inner sep=0.5pt, font=\scriptsize\sffamily},
  lab/.style={font=\scriptsize\sffamily, align=center}]
  % \ptrblk{name}{x}{y}: block of pointers
  \def\ptrblk#1#2#3{%
    \draw[fill=ln-box] (#2-5,#3-4) rectangle (#2+5,#3+4);
    \foreach \k in {-2,0,2} \draw[ln-rule] (#2-5,#3+\k) -- (#2+5,#3+\k);
    \coordinate (#1w) at (#2-5,#3);
    \coordinate (#1e) at (#2+5,#3+3);
    \coordinate (#1s) at (#2+5,#3-3);}
  % inode
  \node[font=\footnotesize\sffamily\bfseries, anchor=south] at (19,48.3) {\iflangja{iノード}{inode}};
  \draw[fill=ln-accent!10] (4,36) rectangle (34,48);
  \node[lab] at (19,42) {\iflangja{メタデータ\\（モード，所有者，\\タイムスタンプ，サイズ）}{metadata\\(mode, owners,\\timestamps, size, \ldots)}};
  \def\irow#1#2{\draw[fill=white] (4,#1-2.5) rectangle (34,#1+2.5);
    \node[lab] at (19,#1) {#2};}
  \irow{33.5}{\iflangja{直接ポインタ 0}{direct 0}}
  \irow{28.5}{\ldots}
  \irow{23.5}{\iflangja{直接ポインタ 11}{direct 11}}
  \irow{18.5}{\iflangja{間接ポインタ}{single indirect}}
  \irow{13.5}{\iflangja{二重間接ポインタ}{double indirect}}
  \irow{8.5}{\iflangja{三重間接ポインタ}{triple indirect}}
  % direct
  \node[data] (a) at (111,33.5) {\iflangja{データ}{data}};
  \node[data] (b) at (111,23.5) {\iflangja{データ}{data}};
  \draw[->] (31,33.5) -- (a.west);
  \draw[->] (31,23.5) -- (b.west);
  % single indirect
  \ptrblk{p}{53}{15}
  \node[data] (c) at (111,15) {\iflangja{データ}{data}};
  \draw[->] (31,18.5) -- (pw);
  \draw[->] (pe) -- (c.west);
  % double indirect
  \ptrblk{q}{53}{0}
  \ptrblk{r}{77}{0}
  \node[data] (d) at (111,0) {\iflangja{データ}{data}};
  \draw[->] (31,13.5) -- (qw);
  \draw[->] (qe) -- (rw);
  \draw[->] (re) -- (d.west);
  % triple indirect
  \ptrblk{s}{53}{-15}
  \ptrblk{t}{71}{-15}
  \ptrblk{u}{89}{-15}
  \node[data] (e) at (111,-15) {\iflangja{データ}{data}};
  \draw[->] (31,8.5) -- (sw);
  \draw[->] (se) -- (tw);
  \draw[->] (te) -- (uw);
  \draw[->] (ue) -- (e.west);
  % legend
  \draw[fill=ln-box] (4,-27) rectangle (10,-23);
  \node[lab, anchor=west] at (11,-25) {\iflangja{ポインタのブロック}{block of pointers}};
  \draw[fill=white] (44,-27) rectangle (50,-23);
  \node[lab, anchor=west] at (51,-25) {\iflangja{データブロック}{data block}};
\end{tikzpicture}
   (width ~116 mm)
\begin{tikzpicture}[x=1mm, y=1mm, font=\scriptsize\sffamily,
  >={Stealth[length=1.4mm]},
  data/.style={draw, fill=white, minimum width=10mm, minimum height=4.5mm,
               inner sep=0.5pt, font=\scriptsize\sffamily},
  lab/.style={font=\scriptsize\sffamily, align=center}]
  % \ptrblk{name}{x}{y}: block of pointers
  \def\ptrblk#1#2#3{%
    \draw[fill=ln-box] (#2-5,#3-4) rectangle (#2+5,#3+4);
    \foreach \k in {-2,0,2} \draw[ln-rule] (#2-5,#3+\k) -- (#2+5,#3+\k);
    \coordinate (#1w) at (#2-5,#3);
    \coordinate (#1e) at (#2+5,#3+3);
    \coordinate (#1s) at (#2+5,#3-3);}
  % inode
  \node[font=\footnotesize\sffamily\bfseries, anchor=south] at (19,48.3) {\iflangja{iノード}{inode}};
  \draw[fill=ln-accent!10] (4,36) rectangle (34,48);
  \node[lab] at (19,42) {\iflangja{メタデータ\\（モード，所有者，\\タイムスタンプ，サイズ）}{metadata\\(mode, owners,\\timestamps, size, \ldots)}};
  \def\irow#1#2{\draw[fill=white] (4,#1-2.5) rectangle (34,#1+2.5);
    \node[lab] at (19,#1) {#2};}
  \irow{33.5}{\iflangja{直接ポインタ 0}{direct 0}}
  \irow{28.5}{\ldots}
  \irow{23.5}{\iflangja{直接ポインタ 11}{direct 11}}
  \irow{18.5}{\iflangja{間接ポインタ}{single indirect}}
  \irow{13.5}{\iflangja{二重間接ポインタ}{double indirect}}
  \irow{8.5}{\iflangja{三重間接ポインタ}{triple indirect}}
  % direct
  \node[data] (a) at (111,33.5) {\iflangja{データ}{data}};
  \node[data] (b) at (111,23.5) {\iflangja{データ}{data}};
  \draw[->] (31,33.5) -- (a.west);
  \draw[->] (31,23.5) -- (b.west);
  % single indirect
  \ptrblk{p}{53}{15}
  \node[data] (c) at (111,15) {\iflangja{データ}{data}};
  \draw[->] (31,18.5) -- (pw);
  \draw[->] (pe) -- (c.west);
  % double indirect
  \ptrblk{q}{53}{0}
  \ptrblk{r}{77}{0}
  \node[data] (d) at (111,0) {\iflangja{データ}{data}};
  \draw[->] (31,13.5) -- (qw);
  \draw[->] (qe) -- (rw);
  \draw[->] (re) -- (d.west);
  % triple indirect
  \ptrblk{s}{53}{-15}
  \ptrblk{t}{71}{-15}
  \ptrblk{u}{89}{-15}
  \node[data] (e) at (111,-15) {\iflangja{データ}{data}};
  \draw[->] (31,8.5) -- (sw);
  \draw[->] (se) -- (tw);
  \draw[->] (te) -- (uw);
  \draw[->] (ue) -- (e.west);
  % legend
  \draw[fill=ln-box] (4,-27) rectangle (10,-23);
  \node[lab, anchor=west] at (11,-25) {\iflangja{ポインタのブロック}{block of pointers}};
  \draw[fill=white] (44,-27) rectangle (50,-23);
  \node[lab, anchor=west] at (51,-25) {\iflangja{データブロック}{data block}};
\end{tikzpicture}
